\section{Projection insertion and pruning}\label{sec:pruning}

The tail-transport lemma assumes regularity of two trajectories.  We first
use it on successively longer suffixes to show that inserting a product
projection preserves regularity.  Two consequences for product starting
states will then allow us to construct regular trajectories by deleting
reflections with small projection probability. Gate deletion by a
telescoping estimate also appears in the Pauli-spectrum approach of
\citet[Lemma~23 and Claim~24]{NadimpalliParhamVasconcelosYuen2024},
where width controls a Frobenius-norm error. Here deletion depends on
projection probability in the actual pruned-prefix state; the retained
layer must then satisfy the regularity condition for every fixed
exponent triple.

\begin{proposition}[Projection insertion]\label{prop:projection-insertion}
Let $W=U_t\cdots U_1$ have a bounded number of layers, and suppose the
trajectory
\[
 \psi_j=U_j\cdots U_1\psi_0,\qquad 0\le j\le t,
\]
is regular, with $\norm{\psi_0}\le1$.  For any product test $B$ at the
output, set
\[
 \chi_t=P_B\psi_t,\qquad
 \chi_j=U_{j+1}^{\dagger}\cdots U_t^{\dagger}\chi_t
 \quad(0\le j<t).
\]
Then the trajectory $(\chi_j)_{j=0}^t$ is regular for $W$.
\end{proposition}

\begin{proof}
All the vectors $\chi_j$ have norm at most one.  We prove their regularity
at the layers in decreasing order of the layer index.  Suppose regularity
has already been established at every layer strictly after layer $j$.
There is nothing to prove if $U_j$ is a layer of one-qubit unitaries.
Otherwise, fix $0\le b<1$ and choose, at each member of the circuit
sequence, a set $S$ of at most $m^b$ reflections in this layer.  Let
$H_A=H_S$ be their joint count and put
\[
 V=U_t\cdots U_{j+1}.
\]
This is the suffix \emph{strictly after} layer $j$, and
$\chi_j=V^{\dagger}P_BV\psi_j$.  The trajectory of $\psi_j$ on $V$ is
regular by hypothesis; that of $\chi_j$ on $V$ is regular by the
induction hypothesis.  Regularity of the original trajectory at layer
$j$ also gives, for every fixed $b<h<1$ and $0<p<h$,
\[
 \norm{\one_{H_A\ge m^h}\psi_j}\le \exp(-m^p)
 \quad\text{eventually}.
\]
Lemma~\ref{lem:tail-transport} therefore gives the same estimate for
$\chi_j$.  The count $H_S$ commutes with every reflection of layer $j$,
so this estimate holds on $\chi_{j-1}$ as well.

The chosen sets $S$ may vary with the circuit sequence.  For each fixed
triple $(b,h,p)$, Lemma~\ref{lem:uniformization} consequently makes the
estimate eventual and uniform over all eligible $S$.  This proves
regularity at layer $j$ and completes the induction.  Its first step
uses the empty suffix, whose regularity hypotheses are vacuous.
\end{proof}

\begin{corollary}[Tail at a product input]\label{cor:product-tail}
Suppose $\psi_0$ is a normalized product state and its trajectory under
$W$ is regular.  Let $P_0=|\psi_0\rangle\langle\psi_0|$ be the full
product projection and $H_0$ its count.  For any product test $B$ at the
output and every fixed $0<p<h<1$,
\begin{equation}\label{eq:product-tail}
 \norm{\one_{H_0\ge m^h}W^{\dagger}P_BW\psi_0}
 \le \exp(-m^p)
 \quad\text{eventually}.
\end{equation}
\end{corollary}

\begin{proof}
The original trajectory is regular by assumption and the projected
trajectory is regular by Proposition~\ref{prop:projection-insertion}.
Since $H_0\psi_0=0$, the input-tail hypothesis of
Lemma~\ref{lem:tail-transport} holds with $b=0$.  Applying that lemma to
$V=W$ proves the claim.
\end{proof}

\begin{corollary}[Overlap and output tail]\label{cor:overlap-tail}
Under the hypotheses of Corollary~\ref{cor:product-tail}, write
$\psi_t=W\psi_0$.  For any product test $B$ at the output and every
fixed $0<p<h<1$,
\begin{equation}\label{eq:overlap-tail}
 \norm{P_B\psi_t}^{2}
 \norm{\one_{H_B\ge m^h}\psi_t}
 \le \exp(-m^p)
 \quad\text{eventually}.
\end{equation}
For each fixed $(h,p)$, both this estimate and
\eqref{eq:product-tail} hold eventually uniformly over the output test
$B$.
\end{corollary}

\begin{proof}
Put $\chi_t=P_B\psi_t$ and $q_B=\norm{P_B\psi_t}^2$.  Apply
Lemma~\ref{lem:tail-transport} to the reverse circuit $V=W^{\dagger}$,
starting at $\chi_t$, with initial count $H_B$ and final projection
$P_0$.  The starting vector has zero $H_B$ count.  Its trajectory is
regular by Proposition~\ref{prop:projection-insertion} and reversal.
The projected-and-returned vector is exactly
\begin{align*}
 W P_0 W^{\dagger}\chi_t
 &=\psi_t\,\ip{\psi_0}{W^{\dagger}P_BW\psi_0}\\
 &=q_B\psi_t.
\end{align*}
Its reverse trajectory is the original trajectory multiplied by
$q_B\in[0,1]$, hence is regular too.  The lemma, with $b=0$, now gives
\eqref{eq:overlap-tail}.  In particular, the multiplier is the
\emph{probability} $q_B$, not its square root: no projected vector has
been normalized.

For uniformity, fix $(h,p)$ and suppose that either estimate fails for
some test at arbitrarily large indices.  Choose a violating test at
those indices and any test at the others.  The resulting single
sequence of tests contradicts the estimate just proved.  Thus tests
may be selected using the exact circuit and its output state.  The
threshold may depend on the fixed exponents; no threshold uniform over
all real exponent pairs is asserted.
\end{proof}

\newpage
\begin{lemma}[A retained layer is regular]\label{lem:retained-layer}
Let $W$ have a regular trajectory from a normalized product state, and
write $\psi=W\psi_0$.  Consider a proposed next layer of disjoint
reflections $I-2P_g$.  Retain precisely the gates satisfying
\begin{equation}\label{eq:retention-rule}
 q_g:=\norm{P_g\psi}^{2}\ge m^{-6}.
\end{equation}
The retained layer is regular on $\psi$, and hence after its action on
$\psi$ as well.
\end{lemma}

\begin{proof}
The mean bound and Markov's inequality put substantial mass on
low-weight product-basis outcomes. Counting these outcomes yields
one of sufficiently large probability; recentering the test there
changes every count by at most the outcome's weight, so
Corollary~\ref{cor:overlap-tail} supplies the required tail estimate.

We first bound the mean count of each retained gate.  Fix $\epsilon>0$
and choose fixed exponents $0<p_0<h_0<\min\{\epsilon,1\}$.  Applied to
the test of a retained gate, \eqref{eq:overlap-tail} gives
\[
 \norm{\one_{H_g\ge m^{h_0}}\psi}
 \le m^6\exp(-m^{p_0}).
\]
Since $0\le H_g\le mI$ and $\norm{\psi}=1$, spectral decomposition
therefore gives
\begin{equation}\label{eq:retained-mean}
 \ip{\psi}{H_g\psi}
 \le m^{h_0}+m^{13}\exp(-2m^{p_0})
 \le m^{\epsilon}
 \quad\text{eventually}.
\end{equation}
Corollary~\ref{cor:overlap-tail} makes this bound uniform over the
retained gates.

Now fix a regularity triple $0\le b<h<1$, $0<p<h$, and choose
$\epsilon>0$ so that $b+2\epsilon<h$.  For any set $S$ of at most
$m^b$ retained gates, disjointness of the supports gives
\[
 H_S=\sum_{g\in S}H_g,
 \qquad
 \ip{\psi}{H_S\psi}\le m^{b+\epsilon}.
\]
The empty set has zero tail, so we may suppose $S$ is nonempty.
Consider the distribution obtained by measuring the tested sites in
the product basis consisting of each specified state and its
orthogonal complement.  Label the outcomes by
binary strings, with $0$ for the specified state.  The Hamming weight
of the resulting string is the measured value of $H_S$.  With
\[
 R=2m^{b+\epsilon},\qquad \beta=b+2\epsilon,
\]
Markov's inequality puts at least half the probability on strings of
weight at most $R$.  If $M\le m$ sites are tested, the number $N_R$ of
such strings satisfies, eventually,
\begin{align*}
 N_R\le\sum_{j=0}^{\lfloor R\rfloor}\binom mj
 &\le (\lfloor R\rfloor+1)m^{\lfloor R\rfloor},\\
 \log(2N_R)&=O(m^{b+\epsilon}\log m)=o(m^{\beta}).
\end{align*}
Consequently one such outcome $z$ has probability
\begin{equation}\label{eq:pruning-atom}
 \rho\ge\frac{1}{2N_R}\ge\exp(-m^{\beta})
 \quad\text{eventually}.
\end{equation}
Only linearity of the mean, Markov's inequality, and counting have
been used; the outcomes need not be independent.

Let $B$ specify the local states of $z$ on the tested sites and leave
all other sites unconstrained.  Then $\rho=\norm{P_B\psi}^2$.
The operators $H_B$ and $H_S$ are diagonal in the same product basis.
At a site where $z$ is $0$ their count terms agree; at a site where
$z$ is $1$ they differ by an operator of norm one.  Thus their
eigenvalues on every common basis vector differ by at most
$|z|\le R$.  Choose fixed exponents
\[
 \max\{p,\beta\}<p'<h'<h.
\]
Because $b+\epsilon<h'<h$, eventually
$m^h-R\ge m^{h'}$.  The common basis then gives the spectral inclusion
\[
 \one_{H_S\ge m^h}\le\one_{H_B\ge m^{h'}}.
\]
Apply \eqref{eq:overlap-tail} to $B$ with exponents $(h',p')$, and
divide by its probability from \eqref{eq:pruning-atom}:
\begin{align*}
 \norm{\one_{H_S\ge m^h}\psi}
 &\le\norm{\one_{H_B\ge m^{h'}}\psi}\\
 &\le\rho^{-1}\exp(-m^{p'})\\
 &\le\exp(m^{\beta}-m^{p'})
 \le\exp(-m^p)
 \quad\text{eventually}.
\end{align*}
The last inequality uses both $p'>\beta$ and $p'>p$.

For each fixed target triple, the preceding bounds are uniform over
$S$ and its selected outcome $z$.  Indeed the endpoint estimate is
uniform over all product tests, including those selected from the
current state.  Equivalently, a sequence of violating choices would
contradict Corollary~\ref{cor:overlap-tail} on that subsequence.  Only
finitely many auxiliary exponent choices were made for this triple.
All uses of the corollary concern the already regular prefix $W$;
the proposed next layer is not among its hypotheses.  We have proved
regularity before the retained layer, and commutation of $H_S$ with
that layer gives the same property afterward.
\end{proof}

\begin{proposition}[Pruning from a product state]\label{prop:pruning}
For any sequence of circuits $U=U_t\cdots U_1$ on $m$ qubits, with a
bounded number $t$ of analytic layers, and any normalized product
starting states $\psi_0$, the following construction produces a
circuit $W$ with a regular trajectory from $\psi_0$.  Leave all
one-qubit layers unchanged.  At a reflection layer, use the exact
state reached by the pruned prefix and retain exactly the gates in
\eqref{eq:retention-rule}.  The output vectors satisfy
\begin{equation}\label{eq:pruning-error}
 \norm{U\psi_0-W\psi_0}\le 2t m^{-2}.
\end{equation}
\end{proposition}

\begin{proof}
The empty prefix is regular.  Appending a one-qubit layer adds no
regularity condition, and Lemma~\ref{lem:retained-layer} applies at
each reflection layer.  Induction therefore gives regularity of the
whole pruned trajectory.  The retention rule itself has no exponent
choices: the auxiliary exponents above are used only to verify its
regularity for each fixed triple.

For the error estimate, denote the retained version of layer $U_j$ by
$\widetilde U_j$ and the exact pruned-prefix state before it by
$\psi_{j-1}$.  Each deleted reflection $F_g=I-2P_g$ satisfies
\[
 \norm{(F_g-I)\psi_{j-1}}
 =2\norm{P_g\psi_{j-1}}<2m^{-3}.
\]
Every other gate in that layer commutes with $P_g$, and is unitary.
Hence this norm is unchanged if any other gates of the layer have
already acted.  Telescoping over the deleted gates, of which there
are at most $m$ because their nonempty supports are disjoint, yields
\[
 \norm{(U_j-\widetilde U_j)\psi_{j-1}}\le2m^{-2}.
\]
This bound is zero for a one-qubit layer.  If $\phi_j$ and $\psi_j$
are the original and pruned states after layer $j$, respectively,
then exactly
\[
 \phi_j-\psi_j
 =U_j(\phi_{j-1}-\psi_{j-1})
   +(U_j-\widetilde U_j)\psi_{j-1}.
\]
Unitarity and the triangle inequality give
$\norm{\phi_j-\psi_j}\le\norm{\phi_{j-1}-\psi_{j-1}}+2m^{-2}$.
Starting with $\phi_0=\psi_0$ proves \eqref{eq:pruning-error}.
\end{proof}
